% ============================================================
% Cấu hình huấn luyện và lượng tử hóa mô hình
%
% Tổng hợp quy trình huấn luyện theo mã nguồn training.py
% ============================================================



\section{Cấu hình huấn luyện và lượng tử hóa mô hình}

Sau khi dữ liệu được tiền xử lý và chuẩn hóa về dải $[0, 1]$, bước tiếp theo là xây dựng quy trình huấn luyện và nén mô hình nhằm đáp ứng các giới hạn về bộ nhớ và năng lực tính toán của Seeed Studio XIAO ESP32-S3.

\subsection{Phân chia tập huấn luyện, kiểm định và kiểm thử}

Quá trình huấn luyện được thực hiện riêng cho từng pha vận hành nhằm tránh việc mô hình bị chi phối bởi trạng thái GENTLE vốn chiếm tỷ trọng lớn trong dữ liệu thô. Từ tập \texttt{train\_features\_v6.csv} gồm 14.000 vectơ đặc trưng 19 chiều, dữ liệu được phân hoạch theo phương sai trục Z thành ba tập GENTLE, STRONG và SPIN như đã trình bày ở Mục~\ref{subsec:data_distribution_balance}. Sau đó, mỗi tập trạng thái được tăng cường độc lập trước khi đưa vào mô hình tự mã hóa tương ứng.

Trong mã huấn luyện, tập kiểm định được tách trực tiếp từ dữ liệu sau tăng cường bằng tham số \texttt{validation\_split=0.15}. Như vậy, 85\% dữ liệu tăng cường được dùng để cập nhật trọng số mô hình, còn 15\% được giữ lại để theo dõi sai số kiểm định và kích hoạt cơ chế dừng sớm. Tập kiểm thử cuối cùng không được dùng trong quá trình tối ưu trọng số; các đánh giá định lượng trên tập kiểm thử độc lập và dữ liệu lỗi giả lập được trình bày ở Chương 4.

\begin{table}[H]
    \centering
    \caption{Cách chia dữ liệu huấn luyện và kiểm định theo từng pha vận hành.}
    \label{tab:train_val_split}
    \begin{tabular}{lcccc}
        \hline
        \textbf{Pha} & \textbf{Mẫu gốc} & \textbf{Sau tăng cường} & \textbf{Huấn luyện 85\%} & \textbf{Kiểm định 15\%} \\
        \hline
        GENTLE & 12.853 & 20.000 & 17.000 & 3.000 \\
        STRONG & 551 & 10.000 & 8.500 & 1.500 \\
        SPIN & 596 & 10.000 & 8.500 & 1.500 \\
        \hline
    \end{tabular}
\end{table}

\subsection{Kỹ thuật tăng cường dữ liệu}
Phương pháp học không giám sát yêu cầu tập dữ liệu huấn luyện bao quát được các biến động tự nhiên của môi trường vận hành bình thường. Hệ thống áp dụng các phương pháp tăng cường dữ liệu trực tiếp trên tín hiệu vật lý: Điều chỉnh rung (Vibration Scaling), Dịch chuyển biên độ âm thanh (Audio Volume Shifting) và Rung động từng trục (Per-axis Jittering) (chi tiết tại Mục~\ref{subsec:data_distribution_balance}).

\subsection{Quá trình huấn luyện mô hình số thực (Float32)}

Mô hình đối xứng được thiết kế theo cấu trúc thắt cổ chai sâu (Deep Bottleneck) với số lượng nơ-ron qua các lớp lần lượt là: $19 \rightarrow 128 \rightarrow 64 \rightarrow 32 \rightarrow 64 \rightarrow 128 \rightarrow 19$. Việc mở rộng không gian đặc trưng ở lớp ẩn đầu tiên (từ 19 lên 128 chiều) tạo cấu trúc biểu diễn dư thừa, cung cấp đủ không gian để trích xuất các mối tương quan phi tuyến giữa tín hiệu âm thanh và rung động trước khi nén tại lớp cổ chai 32 chiều.

\begin{itemize}
    \item \textbf{Hàm kích hoạt Sigmoid:} Mặc dù ReLU là hàm kích hoạt phổ biến, kiến trúc này sử dụng Sigmoid cho toàn bộ các lớp ẩn. Sigmoid giới hạn giá trị đầu ra trong khoảng $[0, 1]$, thu hẹp dải động của dữ liệu và giảm sai số khi thực hiện lượng tử hóa về định dạng số nguyên 8-bit (INT8).
    
    \item \textbf{Hàm mất mát có trọng số:} Trong vector đầu vào, 13 đặc trưng âm thanh chiếm ưu thế về số lượng, nhưng 6 đặc trưng rung động lại là chỉ báo vật lý quan trọng nhất để xác định lỗi cơ khí. Đồ án định nghĩa hàm \texttt{weighted\_mae()}, đặt trọng số hệ số 5 cho các sai số thuộc nhóm gia tốc, nhằm ưu tiên tối ưu hóa khả năng tái tạo động hình cơ học của lồng giặt.
    
    \item \textbf{Thuật toán tối ưu và khử nhiễu:} Mô hình sử dụng thuật toán tối ưu Adam \cite{ref_adam} với tốc độ học khởi tạo $0{,}001$ trong giai đoạn huấn luyện số thực. Để tăng tính bền vững, dữ liệu đầu vào được cộng thêm nhiễu Gaussian nhẹ có độ lệch chuẩn $0{,}02$ trong miền đã chuẩn hóa $[0,1]$, trong khi hàm mục tiêu yêu cầu mô hình tái tạo lại tín hiệu sạch. Cách tiếp cận này giúp mô hình học được các đặc trưng cốt lõi và tránh hội tụ về hàm đồng nhất (identity function).
\end{itemize}

\begin{table}[H]
    \centering
    \caption{Cấu hình huấn luyện được sử dụng trong mã nguồn \texttt{training.py}.}
    \label{tab:training_config}
    \begin{tabular}{lcc}
        \hline
        \textbf{Thông số} & \textbf{Tiền huấn luyện Float32} & \textbf{Tinh chỉnh QAT INT8} \\
        \hline
        Thuật toán tối ưu & Adam & Adam \\
        Tốc độ học & $1\times10^{-3}$ & $1\times10^{-4}$ \\
        Số epoch tối đa & 200 & 80 \\
        Batch size & 64 & 64 \\
        Tỷ lệ kiểm định & 15\% & 15\% \\
        Dừng sớm & Patience 20 & Patience 15 \\
        Khôi phục trọng số tốt nhất & Có & Có \\
        \hline
    \end{tabular}
\end{table}


\subsection{Lượng tử hóa nhận thức huấn luyện}

Để triển khai mô hình trên vi điều khiển Seeed Studio XIAO ESP32-S3, kích thước trọng số phải được nén. Thay vì sử dụng lượng tử hóa sau huấn luyện (PTQ), đồ án áp dụng lượng tử hóa nhận thức huấn luyện. Phương pháp này tích hợp các phép toán mô phỏng lượng tử hóa vào quá trình tinh chỉnh (Fine-tuning), giúp mô hình điều chỉnh trọng số để bù cho sai số do giảm độ phân giải số học.

Quá trình huấn luyện được chia thành hai giai đoạn:
\begin{enumerate}
    \item \textbf{Giai đoạn 1 - Tiền huấn luyện:} Mô hình được huấn luyện với định dạng số thực trong tối đa 200 epochs, batch size 64 và cơ chế dừng sớm sau 20 epochs nếu sai số kiểm định không cải thiện. Sai số kiểm định (Val Loss) hội tụ ở mức xấp xỉ \textbf{0,0210} (đối với pha SPIN).
    \item \textbf{Giai đoạn 2 - Tinh chỉnh:} Mô hình sau tiền huấn luyện được bọc bằng lớp lượng tử hóa nhận thức huấn luyện và tiếp tục tối ưu thêm tối đa 80 epochs với tốc độ học $1\times10^{-4}$. Cơ chế dừng sớm được đặt patience 15. Nhờ QAT, sai số kiểm định chỉ tăng không đáng kể lên mức \textbf{0,0225}.
\end{enumerate}
Mức suy giảm hiệu năng dưới 7\% chứng tỏ QAT bảo toàn độ chính xác của mô hình khi chuyển đổi sang môi trường nhúng.

Sau khi huấn luyện, các mô hình được xuất dưới dạng mảng C (Mảng C). Mỗi mô hình đơn lẻ (25.779 tham số) sau lượng tử hóa INT8 có kích thước xấp xỉ \textbf{31,7 KB}. Tổng dung lượng cho tổ hợp 3 mô hình (hỗn hợp Ba trạng thái) chỉ chiếm \textbf{95,3 KB}, giảm 4 lần so với mô hình Float32, cho phép triển khai thuật toán chuyển đổi trạng thái mà không tăng áp lực lên bộ nhớ Flash.

\subsection{Xác định ngưỡng cảnh báo dựa trên quy tắc $3\sigma$}
Ngưỡng phát hiện dị thường ($\tau$) cho từng pha được thiết lập từ phân phối sai số tái tạo của tập dữ liệu bình thường, tuân theo quy tắc ba sigma ($3\sigma$):
\begin{equation}
    \tau = \mu_{MAE} + 3 \cdot \sigma_{MAE}
\end{equation}

Lựa chọn hệ số $k=3$ thay vì các ngưỡng tối ưu theo chỉ số Youden's J phản ánh yêu cầu thực tiễn: cảnh báo giả (False Positive) gây ảnh hưởng tiêu cực đến người dùng nghiêm trọng hơn phản hồi chậm đối với một số dị thường nhỏ (False Negative). Ngưỡng $3\sigma$ cung cấp dải biên an toàn lớn, đảm bảo độ chính xác (Precision) cao trước khi phát tín hiệu hỏng hóc.

Các ngưỡng cố định được xác lập tương ứng: GENTLE ($0{,}1241$), STRONG ($0{,}0968$), và SPIN ($0{,}0672$).


% ============================================================
% MỤC: XUẤT THÔNG SỐ CẤU HÌNH VÀ THAM SỐ MÔ HÌNH CHO VI ĐIỀU KHIỂN
% ============================================================


\subsection{Đánh giá cấu hình và hiệu năng tài nguyên}

Kiến trúc mạng nơ-ron sử dụng toàn bộ hàm kích hoạt Sigmoid để giới hạn dữ liệu trong dải $[0, 1]$, ngăn ngừa các lỗi tràn bộ đệm khi vi điều khiển thực thi các phép toán nhân ma trận INT8.

\begin{table}[H]
    \centering
    \caption[Chi tiết tham số mạng tự mã hóa]{Chi tiết cấu trúc các lớp và số lượng tham số của một mô hình tự mã hóa thành phần.}
    \label{tab:ae_params}
    \begin{tabular}{clcc}
        \hline
        \textbf{STT} & \textbf{Lớp} & \textbf{Kích thước} & \textbf{Tham số} \\
        \hline
        1 & Đầu vào      & $19$    & ---    \\
        2 & Dense (Encoder-1)    & $128$   & $2.560$ \\
        3 & Dense (Encoder-2)    & $64$    & $8.256$ \\
        4 & Dense (lớp thắt cổ chai)   & $32$    & $2.080$ \\
        5 & Dense (Decoder-1)    & $64$    & $2.112$ \\
        6 & Dense (Decoder-2)    & $128$   & $8.320$ \\
        7 & Đầu ra      & $19$    & $2.451$ \\
        \hline
        \multicolumn{3}{r}{\textbf{Tổng cộng:}} & \textbf{25.779} \\
        \hline
    \end{tabular}
\end{table}

Kích thước tối ưu này (95,3 KB cho 3 mô hình) giúp thời gian tính toán suy luận trên lõi Xtensa LX7 của vi điều khiển Seeed Studio XIAO ESP32-S3 đạt mức khoảng 6,4 ms (Bảng~\ref{tab:latency_stats}), chỉ chiếm 0,64\% quỹ thời gian của một chu kỳ lấy mẫu 1 Hz, duy trì tần số lấy mẫu tốc độ cao mà không làm chậm các tiến trình khác trong hệ điều hành FreeRTOS.


% --- Kết thúc chương 3 mục 4 ---
